\chapter{Меню кандидатов новой модели спектрального рождения}
\label{ch:v6-spectral-transition-new-model-candidate-menu-gate}

\begin{remark}[Ретроспективный статус]
Выбор динамического $D_F(x)$ в этой главе впоследствии отменён гейтом
\path{version6_spectral_transition_candidate_menu_retrospective_correction_gate}:
Том V уже проверил операторнозначный дефект, его нулевую моду,
проекторное сокращение и спектральные источники коэффициентов.
\end{remark}

\section{Правило сравнения}

Контракт R0--R6 не позволяет выбрать модель по одному удачному образу.
Каждый кандидат проверяется по трём вопросам:
\begin{enumerate}
 \item сохраняет ли он физический носитель и уже доказанные классы;
 \item способен ли один функционал дать trigger, скорость и endpoint;
 \item сколько новых полей, зарядов и свободных коэффициентов приходится
 ввести до появления первого нового observable.
\end{enumerate}

Статус \emph{перспективен} ниже означает только отсутствие уже известного
структурного запрета. Он не означает прохождение R0--R6.

\section{Кандидат A: электрослабый сфалерон}

Стандартный gauge--Higgs функционал уже имеет конечное седло и механизм
спектрального потока. Но он действует только на слабые левые дублеты,
не связывает полный пакет $15$ и имеет неустойчивый endpoint. Поэтому
сфалерон остаётся контрольным примером для R1--R3, но окончательно
отклоняется как модель наблюдаемой частицы.

\section{Кандидат B: явный портал между вихрем и Хиггсом}

Можно вручную добавить бифундаментальное поле или член
\begin{equation}
 \lambda_{QH}\Tr(Q^2)H^\dagger H.
\end{equation}
Однако предыдущие гейты доказали отсутствие такого интертвинера и такого
члена в текущем родителе. После ручного добавления коэффициент
$\lambda_{QH}$ становится новым входом; trigger, скорость и численное
предсказание не следуют. Кандидат отклоняется как скрытый портал.

\section{Кандидат C: Q-ball или фермионный мешок}

Нетопологические солитоны могут быть устойчивы при фиксированном
глобальном заряде \cite{candidate-menu-coleman1985,
candidate-menu-friedberg-lee-sirlin1976}. Это реальный механизм конечного
трёхмерного комка, но он требует нового комплексного поля, непрерывного
$U(1)$-заряда и специальной формы потенциала. Проект пока не выводит ни
этот заряд, ни потенциал из KO6/Toeplitz-ledger. Кандидат сохраняется как
контроль R4, но отклоняется как минимальное продолжение.

\section{Кандидат D: фермионная мода на топологическом дефекте}

Массовый оператор, меняющий знак или топологический класс, способен
связывать локализованную фермионную моду; классический пример --- механизм
Jackiw--Rebbi \cite{candidate-menu-jackiw-rebbi1976}. Он естественно
согласуется с языком спектрального потока, но требует сначала построить
сам устойчивый монопольный, струнный или стеночный фон. Текущие проектные
дефекты либо протяжённы, либо не соединены с $\mathcal H_{15}$.
Следовательно, это возможный механизм endpoint внутри будущей модели,
но не самостоятельный родитель.

\section{Кандидат E: динамическое поле конечного оператора Дирака}

Минимальное продолжение состоит в том, чтобы не добавлять отдельную
геометрическую верёвку, а повысить уже существующий конечный оператор до
пространственно-временного поля
\begin{equation}
 \mathfrak D_F:x\longmapsto \mathfrak D_F(x)
 \in\operatorname{End}_{G_{\rm SM},J,\gamma}(\mathcal H_{15}).
 \label{eq:v6-candidate-menu-dynamical-dirac-field}
\end{equation}
Его спектральное действие имеет канонический общий вид
\begin{equation}
 S_{\rm spec}[\mathfrak D_F,A,\Psi]
 =\Tr f(D_A(\mathfrak D_F)/\Lambda)
 +\langle\Psi,D_A(\mathfrak D_F)\Psi\rangle,
 \label{eq:v6-candidate-menu-spectral-action}
\end{equation}
что сохраняет зависимость действия от физического оператора и его
внутренних флуктуаций \cite{candidate-menu-chamseddine-connes1997}.

Этот кандидат единственный из меню одновременно:
\begin{itemize}
 \item начинает с уже типизированного $\mathcal H_{15}$;
 \item делает смену ранга свойством динамического поля, а не внешнего пути;
 \item допускает один функционал для вакуума, седла и endpoint;
 \item может включить Jackiw--Rebbi-подобную локализацию как следствие
 дефекта массового оператора, не как отдельную аксиому.
\end{itemize}
Но R2--R5 ещё не пройдены: кинетический член, потенциал, топология вакуума,
седло, prefactor и устойчивый endpoint должны быть выведены.

\section{Кандидат F: дискретная сеть спектральных переходов}

Дискретный граф переходов соответствует исходной интуиции о минимальном
не стягиваемом размере и избегает фундаментальной гладкой нити. Однако
сейчас он требует заново вывести локальность, лоренцеву кинематику,
$G_{\rm SM}$-представления, Real-структуру и непрерывный предел. Поэтому
кандидат не опровергнут, но отложен как перестройка основания, а не
минимальное продолжение версии VI.

\section{Матрица решения}

\begin{center}
\begin{tabular}{>{\raggedright\arraybackslash}p{0.25\textwidth}
                  >{\raggedright\arraybackslash}p{0.16\textwidth}
                  >{\raggedright\arraybackslash}p{0.19\textwidth}
                  >{\raggedright\arraybackslash}p{0.28\textwidth}}
\toprule
Кандидат & R0 & Главный плюс & Решение \\
\midrule
Сфалерон & частично & готовое седло & отклонить: нет endpoint \\
Явный портал & нет & связывает сектора & отклонить: ручной параметр \\
Q-ball & нет & конечный комок & контроль R4 \\
Jackiw--Rebbi & частично & фермионная мода & подмеханизм endpoint \\
Динамический $D_F$ & да & общий оператор & выбрать для разработки \\
Дискретная сеть & открыто & минимальный размер & отложить \\
\bottomrule
\end{tabular}
\end{center}

\begin{theorem}[Архитектурная развилка]
Ни один кандидат ещё не проходит R0--R6. Единственным минимальным
кандидатом без уже доказанного жёсткого конфликта с физическим носителем
является динамическое поле конечного оператора Дирака. Оно выбирается как
следующий объект вывода, но не объявляется физической моделью до построения
родительского действия и прохождения R2--R5.
\end{theorem}

\begin{nogocriterion}
Выбор динамического $D_F$ не разрешает произвольно назначить ему
кинетический член или потенциал. Если коэффициенты не следуют из
спектрального действия, нормировки предыдущих томов или единственного
нового масштаба с новым предсказанием, кандидат отклоняется.
\end{nogocriterion}

\section{Следующий гейт}

Следующий гейт
\path{version6_spectral_transition_candidate_menu_retrospective_correction_gate}
должен проверить этот выбор по полному ledger Тома V до нового построения.

Нулевая гипотеза: спектральное действие не фиксирует достаточный потенциал
без нескольких новых параметров. Стоп-критерий: если нет единого вакуума,
конечной энергии смены ранга и хотя бы одного параметрически чистого
отношения, кандидат не проходит R1 и R5.

\begin{protocol}
\begin{enumerate}
 \item проверенных архитектур: \textbf{6};
 \item полностью допущенных по R0--R6: \textbf{0};
 \item отклонённых как готовая модель: \textbf{3};
 \item сохранённых как подмеханизм или контроль: \textbf{1};
 \item отложенных фундаментальных перестроек: \textbf{1};
 \item выбранных для следующего вывода: \textbf{1};
 \item выбранный кандидат: \textbf{динамическое поле $D_F(x)$};
 \item статус выбора: \textbf{исследовательский приоритет, не допуск}.
\end{enumerate}
\end{protocol}

Машинный сертификат сохранён в
\path{s2t_v6_spectral_transition_new_model_candidate_menu_gate_results.json}.

\begin{thebibliography}{9}
\bibitem{candidate-menu-chamseddine-connes1997}
A.~H.~Chamseddine, A.~Connes,
\emph{The Spectral Action Principle},
Communications in Mathematical Physics 186 (1997), 731--750;
arXiv:hep-th/9606001.

\bibitem{candidate-menu-jackiw-rebbi1976}
R.~Jackiw, C.~Rebbi,
\emph{Solitons with Fermion Number $1/2$},
Physical Review D 13 (1976), 3398--3409.

\bibitem{candidate-menu-coleman1985}
S.~Coleman, \emph{Q-balls},
Nuclear Physics B 262 (1985), 263--283.

\bibitem{candidate-menu-friedberg-lee-sirlin1976}
R.~Friedberg, T.~D.~Lee, A.~Sirlin,
\emph{A Class of Scalar-Field Soliton Solutions in Three Space Dimensions},
Physical Review D 13 (1976), 2739--2761.
\end{thebibliography}